\section{Metodología}
Para alcanzar los objetivos planteados en el proyecto se propone el siguiente flujo de trabajo:

\begin{enumerate}
    \item Adquirir o desarrollar una base de datos orientada a la segmentación de instancias. En caso de ser necesario, llevar a cabo depuración y arreglos en la base de datos para optimizarla en función de la tarea específica.
    \item Adquirir o desarrollar una base de datos para tareas de seguimiento.
    \item Llevar a cabo el entrenamiento de diversos modelos YOLO para segmentación de instancias, utilizando hiperparámetros de entrenamiento por defecto con algunas variaciones.
    \item Realizar validación de los modelos entrenados en la etapa anterior con el fin de establecer un \textit{benchmark} inicial que permita comparar el desempeño de los distintos modelos en diversas condiciones de entrenamiento y orientar la siguiente etapa.
    \item Ejecutar una búsqueda de hiperparámetros para identificar la configuración óptima para el entrenamiento definitivo del modelo. Dicha búsqueda se restringe a los casos más prometedores identificados en la etapa anterior, con el objetivo de reducir el espacio de búsqueda debido a limitaciones de tiempo.
    \item Proceder al entrenamiento final del modelo, empleando los hiperparámetros optimizados, y validar los resultados obtenidos.
    \item Validar los modelos entrenados en la tarea de seguimiento y evaluar su desempeño teórico.
\end{enumerate}

Cada etapa conlleva desafíos específicos, requiere el uso de herramientas particulares y demanda el desarrollo correspondiente de código. Además, con el objetivo de facilitar la replicación de los experimentos y resultados obtenidos en este trabajo, se contempla de forma paralela el entrenamiento y validación utilizando la base de datos pública Deepfish; para dicha parte no se modifica la base de datos ni se incluyen pruebas relacionadas con seguimiento, a diferencia de lo que ocurre con la base de datos de salmones.